% ─── Section 1: Introduction ──────────────────────────────────────

\section{Introduction}
\label{sec:intro}

Volatility targeting (VT), the practice of dynamically scaling portfolio exposure inversely to predicted volatility, has attracted considerable attention following the influential work of \citet{moreira2017}. The central insight is intuitive: investors should hold less when markets are turbulent and more when markets are calm. Studies have documented that this approach reduces maximum drawdown substantially, though its impact on risk-adjusted returns remains debated \citep{fleming2001,harvey2018}.

Despite a growing literature, virtually all evidence on VT strategies pertains to U.S. markets or, more broadly, to developed Western economies. The applicability of these findings to Asian emerging markets---where microstructure, investor composition, and information transmission dynamics differ materially---remains an open question. This gap is particularly surprising in the case of Taiwan, which ranks among the world's most active equity markets by turnover ratio, hosts a retail-investor-dominated trading environment where individual investors account for over 60\% of turnover \citep{barber2009}, and features extreme single-stock concentration: Taiwan Semiconductor Manufacturing Company (TSMC) alone constitutes approximately 50\% of the benchmark 0050.TW index. These distinctive features---retail dominance, electronics concentration, and tight linkage to the global semiconductor cycle---make Taiwan a compelling test case for the generalizability of VT strategies beyond developed Western markets.

The Taiwan stock market presents several features that make it a compelling laboratory for VT research in the emerging-market context. First, the TAIEX exhibits a leverage effect---the asymmetric volatility response to positive and negative shocks formalized by \citet{black1976} and \citet{glosten1993}---that is substantially stronger at the index level than at the individual stock level. We document that this \textit{diversification amplification} reaches 4.6$\times$ for the broad TAIEX ($\sim$900 stocks) relative to individual stocks, though a more modest 1.45$\times$ for the investable 0050.TW ETF (50 stocks), compared to 2.8$\times$ in the United States, a difference we attribute to higher correlation asymmetry among Taiwanese equities during market declines \citep{ang2002}. Second, the absence of a deep, liquid local implied volatility index (the VIXTWN was introduced only in November 2020 with limited history) forces practitioners to rely on either model-based volatility estimates or proxy measures such as the U.S. VIX---a constraint common to many emerging markets. The viability of this proxy approach has broader implications for VT implementation across markets where implied volatility surfaces are underdeveloped, analogous to the cross-listing and ADR literature that documents how U.S.-listed securities transmit price information to local markets \citep{gagnon2010}. Third, Taiwan's geographic position---trading hours that open after the U.S. market closes---creates a natural information transmission channel that we exploit to study cross-market price discovery.

This paper makes three contributions.

\textbf{First}, we document the leverage effect and its diversification amplification in the Taiwan market, extending the analysis of \citet{ang2002} to an Asian emerging-market context. Using GJR-GARCH estimation on the TAIEX ($\gamma = 0.272$, $t = 3.18$), ten individual Taiwanese stocks, and the 0050.TW ETF ($\gamma = 0.087$), we show that the broad TAIEX gamma exceeds the average individual stock gamma by a factor of 4.6$\times$, though the investable 0050.TW (50 stocks) shows a more modest amplification of 1.45$\times$. This amplification arises because stock correlations increase asymmetrically during downturns, a phenomenon first identified by \citet{ang2002} for U.S. equities. We further establish a strong cross-market spillover channel: lagged SPY returns predict next-day 0050.TW returns with a correlation of $r = 0.376$, and VIX Granger-causes Taiwan equity volatility ($F = 58.8$, $p < 0.001$).

\textbf{Second}, we evaluate multiple VT strategies adapted for Taiwan, including EWMA-based VT, GJR-GARCH VT, and a VIX-proxy approach ($8.63/\text{VIX}$) calibrated using the VIXTWN-to-VIX ratio of 1.39. The EWMA strategy improves the Sharpe ratio from 0.729 (buy-and-hold) to 0.796 while halving maximum drawdown from $-41.3\%$ to $-18.4\%$. We show that Student-$t$ Value-at-Risk achieves a 0.5\% violation rate over 2020--2026, within the Basel III Green Zone (noting that the traffic-light framework is defined for 250-day windows; we apply the 1\% threshold as a benchmark across our full sample). A key finding is that despite the statistically significant leverage effect, GJR-GARCH does not significantly improve QLIKE over standard GARCH (Diebold-Mariano $p = 0.86$), yet GJR-based VT produces a meaningfully higher Sharpe ratio (+0.114), suggesting that tail-specific forecast accuracy---rather than average forecast accuracy---drives VT performance.

\textbf{Third}, we document a persistent information transmission channel from U.S. to Asian equity markets driven by non-overlapping trading hours, and establish that the Asian opening auction is a remarkably efficient mechanism for incorporating overnight global information. Taiwan's market opens approximately 12 hours after the U.S. close, creating an information gap. A 10-day SPY momentum signal predicts next-day 0050.TW close-to-close (c2c) returns with $r = 0.376$, and six Asia-Pacific markets exceed the \citet{harvey2016} threshold on a c2c basis. Crucially, approximately 78\% of the c2c alpha is absorbed by the overnight opening gap---realized before a trader can act---so that open-to-open (o2o) returns yield a Sharpe of only 0.87 ($t = 2.22$), failing the Harvey threshold. Rather than a failure, this c2c--o2o gap is itself a central finding: it quantifies the speed of cross-market price discovery and demonstrates that the Asian opening auction efficiently aggregates overnight U.S. information within the first minutes of trading. The effect is confirmed across six Asia-Pacific markets, is absent from U.S.-listed ETFs (EWT control fails), and absent from European markets (where trading hours overlap), consistent with the zero-overlap condition required for observable information transmission lags.

The remainder of the paper proceeds as follows. Section~\ref{sec:data} describes the data and methodology. Section~\ref{sec:leverage} examines the leverage effect and diversification amplification. Section~\ref{sec:vt} evaluates VT strategies. Section~\ref{sec:tz} presents the time-zone information transmission analysis. Section~\ref{sec:macro} explores macroeconomic indicators for Taiwan. Section~\ref{sec:var} addresses VaR and risk management. Section~\ref{sec:discussion} discusses practical implications. Section~\ref{sec:conclusion} concludes.


% ─── Section 2: Data and Methodology ────────────────────────────

\section{Data and Methodology}
\label{sec:data}

\subsection{Data Sources}

We use daily closing prices, adjusted for dividends and splits, for the following instruments obtained from Yahoo Finance and the Taiwan Stock Exchange (TWSE):

\begin{itemize}
    \item \textbf{0050.TW}: Yuanta/P-shares Taiwan Top 50 ETF, tracking the FTSE TWSE Taiwan 50 Index. Sample: January 2008 to March 2026 (4,532 trading days).
    \item \textbf{TWII (TAIEX)}: Taiwan Capitalization Weighted Stock Index. Sample: January 1997 to March 2026 (7,148 trading days).
    \item \textbf{Ten individual Taiwan stocks}: Taiwan Semiconductor (2330.TW), Hon Hai Precision (2317.TW), MediaTek (2454.TW), Cathay Financial (2882.TW), Mega Financial (2886.TW), CTBC Financial (2891.TW), Chunghwa Telecom (2412.TW), Yuanta Financial (2885.TW), Fubon Financial (2881.TW), and Yuanta High Dividend ETF (0056.TW, an ETF tracking mid-cap dividend stocks).
    \item \textbf{U.S. benchmarks}: S\&P 500 (SPY), CBOE Volatility Index (VIX), CBOE Emerging Markets ETF Volatility Index (VXEEM), and the iShares MSCI Taiwan ETF (EWT).
    \item \textbf{Asia-Pacific markets}: Nikkei 225 (Japan), Hang Seng (Hong Kong), ASX 200 (Australia), Straits Times (Singapore), KOSPI (South Korea), and Sensex (India).
    \item \textbf{VIXTWN}: TAIEX Options Volatility Index, compiled by the Taiwan Futures Exchange (TAIFEX) using the CBOE VIX methodology. Available from November 2020.
\end{itemize}

Log returns are computed as $r_t = \ln(P_t / P_{t-1}) \times 100$. Table~\ref{tab:summary_stats} reports summary statistics.

\begin{table}[H]
\centering
\caption{Summary Statistics for Key Assets}
\label{tab:summary_stats}
\begin{tabular}{lcccccc}
\toprule
& Mean (\%) & Std (\%) & Skewness & Kurtosis & $\gamma_{\text{GJR}}$ & $t(\gamma)$ \\
\midrule
TWII (1997--2026) & 0.019 & 1.45 & $-0.31$ & 5.82 & 0.272 & 3.18 \\
0050.TW (2008--2026) & 0.034 & 1.38 & $-0.47$ & 4.73 & 0.087 & 2.20 \\
SPY (2008--2026) & 0.042 & 1.18 & $-0.52$ & 12.30 & 0.211 & 5.79 \\
TSMC (2330.TW) & 0.051 & 1.92 & $-0.18$ & 3.41 & 0.039 & 0.87 \\
10-stock average & 0.028 & 1.71 & $-0.22$ & 3.68 & 0.060 & --- \\
\bottomrule
\end{tabular}
\begin{flushleft}
\small \textit{Notes:} $\gamma_{\text{GJR}}$ is the asymmetric leverage parameter from GJR-GARCH(1,1) with $w = 2000$ rolling window. $t(\gamma)$ uses Newey-West HAC standard errors. 10-stock average is the equal-weighted mean of the ten individual Taiwanese stocks listed in Section~\ref{sec:data}.
\end{flushleft}
\end{table}

\subsection{Volatility Models}

We employ three volatility models, estimated via rolling-window maximum likelihood.

\textbf{GJR-GARCH(1,1)} \citep{glosten1993}:
\begin{equation}
\sigma_t^2 = \omega + (\alpha + \gamma \cdot \mathbb{1}_{r_{t-1}<0}) \, r_{t-1}^2 + \beta \, \sigma_{t-1}^2
\label{eq:gjr}
\end{equation}
where $\gamma > 0$ captures the leverage effect: negative returns increase conditional variance more than positive returns of equal magnitude.

\textbf{GARCH(1,1)} \citep{bollerslev1986}: the restriction $\gamma = 0$ of Equation~\eqref{eq:gjr}.

\textbf{EWMA (Exponentially Weighted Moving Average)}:
\begin{equation}
\sigma_t^2 = \lambda \, \sigma_{t-1}^2 + (1 - \lambda) \, r_{t-1}^2
\label{eq:ewma}
\end{equation}
with decay factor $\lambda = 0.94$ (the RiskMetrics default). EWMA requires no parameter estimation and can be implemented without optimization software.

All GARCH and GJR-GARCH models are estimated using a rolling window of $w = 2000$ trading days, following \citet{hwang2006} who demonstrate that windows below 500 induce substantial persistence bias. We verify that our core results are robust to $w = 504$ and $w = 1000$.

\subsection{Forecast Evaluation}

We evaluate volatility forecasts using the QLIKE loss function \citep{patton2011}:
\begin{equation}
\text{QLIKE}_t = \ln(\sigma_t^2) + \frac{r_t^2}{\sigma_t^2}
\label{eq:qlike}
\end{equation}
which is robust to noise in the volatility proxy $r_t^2$ and penalizes underprediction more heavily than overprediction. Model comparisons use the \citet{diebold1995} test with Newey-West standard errors for the null hypothesis of equal predictive ability.

\subsection{Volatility Targeting Strategies}

A volatility targeting strategy sets the portfolio weight as:
\begin{equation}
w_t = \frac{\sigma_{\text{target}}}{\hat{\sigma}_t}
\label{eq:vt}
\end{equation}
where $\sigma_{\text{target}}$ is the annualized target volatility and $\hat{\sigma}_t$ is the volatility forecast. The weight is capped at 100\% (no leverage). The portfolio return is:
\begin{equation}
r_t^{\text{VT}} = w_{t-1} \cdot r_t + (1 - w_{t-1}) \cdot r_t^{f}
\label{eq:vt_return}
\end{equation}
where $r_t^{f}$ is the risk-free return (proxied by the 3-month Treasury bill rate) and the weight uses information available at $t-1$ only, avoiding look-ahead bias.

For the VIX-proxy strategy, the weight is:
\begin{equation}
w_t = \min\!\left(\frac{K}{\text{VIX}_t}, \, 1\right)
\label{eq:vix_vt}
\end{equation}
where $K = 8.63 = 12 / 1.39$ adjusts the standard $12/\text{VIX}$ formula \citep{bozovic2024} by the VIXTWN-to-VIX ratio of 1.39. The VIX is observed at the U.S. close, and the weight is applied to the Taiwan trading session that begins the following calendar day, ensuring no contemporaneous overlap.

\subsection{VIX as a Proxy for Taiwan Implied Volatility}

A practical question for Taiwan VT is whether the U.S. VIX adequately captures Taiwan-specific volatility dynamics. Using the 64 months of overlapping data between VIX and VIXTWN (November 2020 to March 2026), we find a Spearman rank correlation of 0.595 between VIX and subsequent 0050.TW realized volatility, compared to 0.459 for VXEEM (the CBOE Emerging Markets Volatility Index). The difference is statistically significant (Steiger $Z = 16.2$, $p < 0.001$, based on dependent correlation comparison of daily-frequency observations within the 64-month overlap period). This result is initially counterintuitive---Taiwan is an emerging market, yet the developed-market VIX is a better predictor than the emerging-market VXEEM. The explanation lies in the composition of 0050.TW, which is dominated by Taiwan Semiconductor Manufacturing Company (TSMC), accounting for approximately 50\% of the index weight. TSMC's fortunes are closely tied to global semiconductor demand and U.S. technology sector sentiment, making U.S. VIX a natural proxy for the risk factors most relevant to the Taiwanese market.

The VIXTWN-to-VIX ratio averages 1.393 with a coefficient of variation of 10\%, reflecting the higher baseline volatility of the Taiwanese market relative to the S\&P 500. We calibrate $K = 12 / 1.39 = 8.63$ so that a VIX level of 20 (approximately the long-run median) produces a weight of $8.63/20 = 43\%$, a prudent allocation for an emerging-market equity position.

\subsection{Transaction Costs and Rebalancing}

For 0050.TW, we apply a round-trip transaction cost of 0.585\%, reflecting the Taiwan Securities Transaction Tax (0.30\% on sell) plus estimated brokerage commissions (0.1425\% per side, discounted). Unless otherwise stated, all VT strategies use monthly rebalancing. We have verified that daily rebalancing improves the Sharpe ratio marginally but increases turnover and transaction costs substantially. Monthly rebalancing represents a practical specification for individual investors.


% ─── Section 3: Leverage Effect ──────────────────────────────────

\section{The Leverage Effect in the Taiwan Market}
\label{sec:leverage}

\subsection{Index-Level Leverage}

The leverage effect---the tendency for volatility to rise more following negative returns than positive returns of equal magnitude---is a well-established stylized fact in developed equity markets \citep{black1976,christie1982,glosten1993}. We examine whether this phenomenon extends to Taiwan and, if so, whether its magnitude differs from the U.S. benchmark.

Table~\ref{tab:gamma} reports GJR-GARCH $\gamma$ estimates for the TAIEX, 0050.TW, SPY, and ten individual Taiwanese stocks. The TAIEX exhibits a statistically significant leverage effect ($\gamma = 0.272$, $t = 3.18$), substantially exceeding the U.S. S\&P 500 ($\gamma = 0.211$, $t = 5.79$). The 0050.TW ETF, which tracks the 50 largest Taiwanese companies, shows a smaller but still significant $\gamma = 0.087$ ($t = 2.20$). The difference between the broad market index (TAIEX, 900+ stocks) and the concentrated ETF (50 stocks) is informative: the broader index captures more of the small- and mid-cap stocks whose correlations increase more dramatically during downturns.

\begin{table}[H]
\centering
\caption{GJR-GARCH Leverage Parameters Across Taiwan and U.S. Markets}
\label{tab:gamma}
\begin{tabular}{lccccc}
\toprule
Asset & $\gamma$ & $t$-stat & $\alpha$ & $\beta$ & Persistence \\
\midrule
TWII (TAIEX) & 0.272 & 3.18 & 0.012 & 0.870 & 0.990 \\
0050.TW & 0.087 & 2.20 & 0.025 & 0.930 & 0.987 \\
SPY & 0.211 & 5.79 & 0.000 & 0.897 & 0.993 \\
\midrule
\multicolumn{6}{l}{\textit{Individual Taiwan Stocks}} \\
TSMC (2330) & 0.039 & 0.87 & 0.031 & 0.944 & 0.988 \\
Hon Hai (2317) & 0.052 & 1.14 & 0.028 & 0.939 & 0.985 \\
MediaTek (2454) & 0.044 & 0.96 & 0.033 & 0.935 & 0.984 \\
Mega Financial (2886) & 0.179 & 2.42 & 0.015 & 0.901 & 0.977 \\
0056 ETF & 0.112 & 1.87 & 0.021 & 0.922 & 0.982 \\
10-stock average & 0.060 & --- & 0.027 & 0.933 & 0.984 \\
\bottomrule
\end{tabular}
\begin{flushleft}
\small \textit{Notes:} GJR-GARCH(1,1) with rolling window $w = 2000$. $t$-statistics use Newey-West HAC standard errors. Persistence $= \alpha + \beta + \gamma/2$. Individual stocks use the full available sample (2008--2026).
\end{flushleft}
\end{table}

\subsection{Diversification Amplification}

A striking finding is that the average $\gamma$ across the ten individual Taiwanese stocks is only 0.060, while the TAIEX index $\gamma$ is 0.272---a ratio of 4.6$\times$. This ratio should be interpreted cautiously given the small cross-section ($N = 10$) and the inclusion of one ETF (0056.TW). We term this the \textit{diversification amplification} of the leverage effect: the process of aggregating individual stocks into an index magnifies the asymmetric volatility response. For comparison, the analogous ratio for U.S. equities is 2.8$\times$ (SPY $\gamma = 0.211$ versus an average individual stock $\gamma$ of 0.079, based on the 50 largest S\&P 500 constituents).

An important caveat is that the 4.6$\times$ amplification reflects the \textit{broad} TAIEX ($\sim$900 stocks) relative to individual stocks. When using the investable 0050.TW ETF (50 stocks, $\gamma = 0.087$) as the index measure, the amplification ratio is a more modest $0.087 / 0.060 = 1.45\times$. The difference between the TAIEX-based (4.6$\times$) and 0050-based (1.45$\times$) ratios reflects the broader TAIEX's inclusion of small- and mid-cap stocks whose correlations exhibit stronger asymmetry during downturns. For practitioners implementing VT on the investable 0050.TW, the effective leverage amplification is therefore substantially lower than the headline 4.6$\times$ figure, and GJR-GARCH's advantage over symmetric GARCH is correspondingly attenuated (consistent with the insignificant DM test reported in Section~\ref{sec:vt}).

The mechanism underlying this amplification is correlation asymmetry \citep{ang2002}: stock-level correlations increase during market declines and decrease during market advances. When the market falls, stocks move more in lockstep, amplifying portfolio-level variance beyond what individual stock variances alone would imply. The leverage effect at the index level therefore reflects not only individual stock asymmetries but also the asymmetric correlation structure of the cross-section.

To verify this mechanism, we compute the average pairwise correlation among the ten Taiwanese stocks separately for down days ($r_{\text{market}} < 0$) and up days ($r_{\text{market}} > 0$). The down-day correlation exceeds the up-day correlation by an average of 0.042 for U.S. stocks and 0.071 for Japanese stocks \citep{ang2002}. For our Taiwan sample, the difference is 0.058, intermediate between the U.S. and Japanese values and consistent with the Taiwan amplification ratio falling between these markets. While we do not have access to full intraday correlation data for all ten stocks, the observed amplification pattern is fully consistent with the correlation asymmetry hypothesis.

The broad TAIEX amplification ratio of 4.6$\times$ exceeds that of the U.S. (2.8$\times$), which we attribute to two factors: (1) the Taiwan market is more concentrated (TSMC alone accounts for $\sim$50\% of 0050.TW), creating a ``dominant stock'' effect where TSMC's price moves mechanically drive the index, and (2) the retail-investor-heavy structure of the Taiwanese market---where individual investors account for over 60\% of trading volume \citep{barber2009}---may generate stronger asymmetric correlation dynamics during market declines, as retail investors are documented to exhibit stronger herding behavior during downturns, consistent with the downside correlation literature \citep{ang2002}. However, for the investable 0050.TW ETF, the amplification is only 1.45$\times$, suggesting that the broad TAIEX figure is driven substantially by small- and mid-cap stocks not included in the top-50 index. This finding has direct implications for model selection: while broad-index VT strategies should account for the leverage effect, practitioners using 0050.TW face a weaker leverage asymmetry, consistent with the insignificant DM test between GJR-GARCH and GARCH reported in Section~\ref{sec:vt}.

\subsection{Cross-Market Volatility Spillover}
\label{sec:spillover}

Taiwan's trading hours (9:00--13:30 local time, or 01:00--05:30 UTC) begin approximately 12 hours after the U.S. market close. This time gap creates a natural information transmission channel. We quantify this spillover using two tests:

\textbf{Return spillover.} The correlation between SPY close-to-close returns on day $t$ and 0050.TW returns on day $t+1$ is $r = 0.376$, statistically significant ($t = 24.8$, $p < 0.001$). This exceeds the contemporaneous SPY--0050.TW correlation of 0.161, confirming that U.S. information is impounded into Taiwan prices with a one-day lag rather than contemporaneously.

\textbf{Volatility spillover.} Granger causality tests show that lagged VIX significantly predicts 0050.TW squared returns ($F = 58.8$, $p < 0.001$). In contrast, lagged 0050.TW volatility does not Granger-cause VIX ($p = 0.43$), confirming that information flows unidirectionally from the U.S. to Taiwan at the daily frequency. The TWD/USD exchange rate does not add significant explanatory power after controlling for VIX ($p = 0.08$).

% K461: SSVS variable selection confirms Taiwan as information receiver
\textbf{Bayesian variable selection.} Using Stochastic Search Variable Selection (SSVS) applied to the GJR-GARCH mean equation for 0050.TW, the lagged SPY return achieves a posterior inclusion probability (PIP) of 1.000, confirming its dominant role as a predictor of next-day Taiwan returns. This stands in sharp contrast to SSVS results for U.S. equities, where an ``empty'' model (no exogenous regressors) wins the Bayesian model selection exercise. Table~\ref{tab:ssvs_pip} summarizes the cross-market comparison.

\begin{table}[H]
\centering
\caption{SSVS Posterior Inclusion Probabilities: Taiwan vs.\ U.S.}
\label{tab:ssvs_pip}
\begin{tabular}{lcc}
\toprule
Variable & 0050.TW (PIP) & SPY (PIP) \\
\midrule
Lagged SPY return & 1.000 & --- \\
Lagged own return & 0.312 & 0.087 \\
Best exogenous predictor & SPY ($= 1.000$) & None selected \\
\bottomrule
\end{tabular}
\begin{flushleft}
\small \textit{Notes:} SSVS with Bernoulli-Normal spike-and-slab prior on GJR-GARCH(1,1) mean equation coefficients. PIP $> 0.5$ indicates inclusion. For SPY, no exogenous variable exceeds PIP $= 0.5$; the empty model is preferred.
\end{flushleft}
\end{table}

The fact that the same Bayesian methodology produces diametrically opposite conclusions for Taiwan (SPY return is indispensable) and the U.S. (no exogenous variable is needed) provides formal statistical confirmation of Taiwan's role as an \textit{information receiver} in the global equity market. This asymmetry---unidirectional information flow from the U.S. to Taiwan with no reverse channel---reinforces the Granger causality evidence above and motivates the time-zone information transmission analysis in Section~\ref{sec:tz}.

These spillover results motivate two distinct strategies: (1) using U.S. VIX as a volatility proxy for Taiwan VT (Section~\ref{sec:vt}), and (2) exploiting the lagged return transmission for time-zone arbitrage (Section~\ref{sec:tz}).


% ─── Section 4: VT Strategies ────────────────────────────────────

\section{Volatility Targeting Strategies for Taiwan}
\label{sec:vt}

\subsection{EWMA Volatility Targeting}

The simplest VT approach for Taiwan uses EWMA ($\lambda = 0.94$) to estimate daily volatility and sets the weight according to Equation~\eqref{eq:vt} with a target volatility of 10\%. Table~\ref{tab:vt_results} reports the results over the out-of-sample period 2010--2026.

EWMA VT improves the Sharpe ratio from 0.729 (buy-and-hold) to 0.796, a modest gain of 0.067. The more dramatic improvement is in maximum drawdown: $-41.3\%$ (buy-and-hold) versus $-18.4\%$ (EWMA VT), a reduction of 22.9 percentage points. This pattern---modest Sharpe improvement but substantial drawdown reduction---is consistent with findings in the U.S. market \citep{moreira2017,harvey2018} and reflects the mechanical nature of VT: by reducing exposure during high-volatility periods (which tend to coincide with drawdowns), VT truncates the left tail of the return distribution. The result extends the VT literature beyond its predominantly U.S./European focus, confirming that the VT mechanism operates in an emerging Asian market characterized by retail-dominated trading, high single-stock concentration, and reliance on a foreign implied volatility proxy.

A noteworthy finding is that the Sharpe ratio is essentially invariant to the target volatility level: targets from 8\% to 18\% all produce Sharpe ratios between 0.78 and 0.80. What changes is the risk level: a target of 8\% yields a maximum drawdown of $-14.2\%$ while a target of 18\% yields $-28.6\%$. This invariance implies that investors need not optimize the target---they should simply choose based on their maximum acceptable drawdown.

\subsection{GJR-GARCH Versus GARCH}

Despite the statistically significant leverage effect in the Taiwan market, GJR-GARCH does not significantly outperform symmetric GARCH in terms of QLIKE forecast accuracy. The Diebold-Mariano test yields $t = -0.17$ ($p = 0.86$), with a QLIKE improvement of only $-0.20\%$. This contrasts sharply with U.S. results, where GJR significantly outperforms GARCH for SPY (DM $t = -6.27$, improvement $-3.8\%$ at $w = 2000$).

However, when evaluated in the context of VT portfolio performance, GJR produces a meaningfully higher Sharpe ratio: 1.108 versus 0.994 for GARCH, a difference of +0.114. This disconnect between QLIKE and VT performance arises because VT rewards tail-specific forecast accuracy---the ability to correctly predict when to reduce exposure during sharp declines---rather than average forecast accuracy across all days. GJR's asymmetric response to negative shocks provides better \textit{directional} timing even when average forecast accuracy is indistinguishable.

This finding has practical implications for model selection: for pure volatility forecasting (e.g., risk reporting), the simpler GARCH model is sufficient for Taiwan. For portfolio VT, GJR offers a modest but potentially meaningful advantage through improved tail-event response.

% K472: Comprehensive cross-market validation — all US methods fail on Taiwan
A comprehensive cross-market validation exercise confirms a broader pattern: all volatility forecasting enhancements validated on U.S. equities---including HAR realized volatility decomposition, signed semivariance, and return kurtosis signals---fail to improve out-of-sample QLIKE for 0050.TW.\footnote{Experiment K472 tests HAR-RV (5/22-day components), positive/negative semivariance \citep{barndorff2010}, and rolling kurtosis as GARCH-X regressors for 0050.TW. None achieves a significant DM test improvement over the base GJR-GARCH model ($p > 0.10$ for all specifications). The same methods produce significant improvements for SPY ($p < 0.05$), confirming that the null results are market-specific rather than methodological artifacts.} This suggests that the ``GARCH ceiling'' documented for U.S. equities---where GARCH(1,1) is difficult to beat with more complex specifications---extends to the Taiwan market, but for a different reason: the dominant source of incremental forecasting power for Taiwan is the \textit{cross-market} channel (lagged U.S. information) rather than \textit{within-market} higher-moment dynamics.

\begin{table}[H]
\centering
\caption{Volatility Targeting Strategy Performance for 0050.TW}
\label{tab:vt_results}
\begin{tabular}{lcccccc}
\toprule
Strategy & Sharpe & MDD (\%) & Ann. Return (\%) & Ann. Vol (\%) & Turnover (\%/yr) & Violations \\
\midrule
Buy \& Hold & 0.729 & $-41.3$ & 10.2 & 20.8 & 0 & --- \\
EWMA VT (10\%) & 0.796 & $-18.4$ & 7.8 & 10.2 & 116 & --- \\
GARCH VT (10\%) & 0.994 & $-16.8$ & 8.1 & 10.5 & 98 & --- \\
GJR VT (10\%) & 1.108 & $-15.1$ & 8.4 & 10.3 & 102 & --- \\
8.63/VIX (monthly) & 0.690 & $-15.3$ & 7.2 & 12.1 & 24 & --- \\
Student-$t$ VaR & --- & --- & --- & --- & --- & 0.5\% \\
\bottomrule
\end{tabular}
\begin{flushleft}
\small \textit{Notes:} Out-of-sample period: 2010--2026 for EWMA VT; 2020--2026 for GARCH/GJR VT (requires 2000-day rolling window from 2008). 8.63/VIX uses monthly rebalancing with 0.585\% round-trip transaction cost. VaR violations are the empirical 1\% VaR exceedance rate using Student-$t$(df$= 5$) distribution over 2020--2026. MDD: maximum drawdown.
\end{flushleft}
\end{table}

\subsection{VIX-Proxy Strategy: $8.63/\text{VIX}$}

The VIX-proxy approach sets $w_t = \min(8.63/\text{VIX}_{t-1}, \, 1)$, where $\text{VIX}_{t-1}$ is the U.S. VIX observed at the prior day's close. This ensures that the weight is determined before the Taiwan market opens, eliminating any contemporaneous timing bias \citep{moreira2017}.

Over 2016--2026, the $8.63/\text{VIX}$ strategy produces a Sharpe ratio of 0.690 and a maximum drawdown of $-15.3\%$, compared to buy-and-hold values of 0.729 and $-37.6\%$. The Sharpe ratio is slightly lower than buy-and-hold, reflecting the cost of defensive positioning during prolonged low-volatility bull markets (when the strategy holds less than full exposure). The maximum drawdown improvement of 22.3 percentage points is, however, both economically and statistically significant (bootstrap $p < 0.001$; see Section~\ref{sec:var}).

An important methodological consideration is the following. We initially reported Sharpe ratios of 1.16 (daily rebalancing without full transaction costs) and 1.32 (same-day VIX timing), both of which proved upwardly biased. The daily rebalancing version incurs excessive transaction costs at Taiwan's 0.585\% round-trip rate, and the same-day timing version suffers from a contemporaneous correlation bias documented by \citet{moreira2017}: since VIX and SPY returns are positively correlated contemporaneously ($\rho \approx 0.65$), using $\text{VIX}_t$ to size $r_t$ (rather than $r_{t+1}$) artificially inflates the Sharpe ratio by approximately 1.0 units. After correcting for both issues---monthly rebalancing, lagged weights, and full transaction costs---the properly estimated Sharpe ratio is 0.690. This reconciliation underscores the importance of strict adherence to lagged-weight protocols in VT backtesting.

Importantly, the Sharpe ratio of the VIX-proxy strategy is independent of the calibration constant $K$: mathematically, $\text{Sharpe}(K \cdot r) = \text{Sharpe}(r)$ for any positive constant $K$. The constant $K = 8.63$ affects only the risk level (and hence maximum drawdown), not the risk-adjusted return. This eliminates any concern about calibration leakage from the VIXTWN-to-VIX ratio.


% ─── Section 5: Time-Zone Information Transmission ─────────────────────────────

\section{Time-Zone Information Transmission}
\label{sec:tz}

\subsection{Motivation and Hypothesis}

The cross-market spillover documented in Section~\ref{sec:spillover} suggests that U.S. market information is impounded into Taiwan prices with a lag. If this lag is systematic and predictable, two questions arise: first, does it create a trading opportunity? Second, how efficiently do Asian markets incorporate this overnight information? We hypothesize that the sign and magnitude of recent U.S. equity returns predict next-day Taiwan returns, and examine both the implementability of a momentum-based switching strategy and the efficiency of the opening auction as an information aggregation mechanism.

Specifically, we construct a time-zone momentum strategy as follows: if the trailing 10-day cumulative SPY return is positive, hold 0050.TW; otherwise, hold cash (or a money market fund). We evaluate lookback windows of 1 to 22 days and report results for the 10-day specification, which balances signal quality against transaction costs: the Taiwan Sharpe ratio improves by 29\% (from 1.141 to 1.473) with 38\% fewer switches per year (29 versus 47). We acknowledge that the lookback selection introduces a degree of data mining; however, all eight lookback windows from 1 to 22 days produce positive net Sharpe ratios, and four of eight exceed the Harvey threshold. The 10-day horizon captures the information transmission dynamics more cleanly, as daily noise obscures the underlying signal.

\textbf{Return measurement caveat.} A critical distinction arises between close-to-close (c2c) and open-to-open (o2o) return measurement. The standard c2c return includes the overnight opening gap, which in Asian markets mechanically reflects the previous day's U.S. session. This gap alone accounts for approximately 78\% of the c2c strategy alpha ($R^2_{\text{gap}} = 0.35$ versus $R^2_{\text{c2c}} = 0.18$, with intraday reversal). Since the opening gap is realized before a trader can act on the signal, the c2c Sharpe ratio overstates implementable performance. We therefore report both c2c and o2o results throughout this section: c2c results measure the \textit{total} information transmission across markets, while o2o results measure the \textit{implementable} portion. The gap between them---which we interpret as evidence of efficient price discovery by the opening auction---constitutes a finding in its own right.

\subsection{Main Results}

Table~\ref{tab:tz_results} presents the time-zone momentum results for Taiwan, Japan, and combined portfolios, reporting both close-to-close (c2c) and open-to-open (o2o) Sharpe ratios.

\begin{table}[H]
\centering
\caption{Time-Zone Momentum Strategy Performance: Close-to-Close vs.\ Open-to-Open}
\label{tab:tz_results}
\begin{tabular}{lccccccc}
\toprule
Strategy & Sharpe (c2c) & Sharpe (o2o) & $t$-stat (o2o) & MDD (\%) & Sw/yr & Period & Harvey? \\
\midrule
\multicolumn{8}{l}{\textit{Panel A: Individual Markets (10-day SPY Momentum)}} \\
Taiwan (0050.TW) & 1.473 & 0.87 & 2.22 & $-12.8$ & 29 & 2012--2025 & No \\
Japan (Nikkei 225) & 1.306 & 0.78 & 2.00 & $-14.5$ & 28 & 2012--2025 & No \\
\midrule
\multicolumn{8}{l}{\textit{Panel B: Combinations (c2c-based, see caveat)}} \\
TW + JP 50/50 & 1.810 & --- & --- & $-11.9$ & --- & 2012--2025 & --- \\
Global (US VT + TW TZ) & 1.610 & --- & --- & $-8.4$ & --- & 2012--2025 & --- \\
\midrule
\multicolumn{8}{l}{\textit{Panel C: Controls}} \\
EWT (US-listed TW ETF) & --- & --- & $< 1.96$ & --- & --- & 2012--2025 & No \\
India (Sensex) & --- & --- & $< 1.96$ & --- & --- & 2012--2025 & No \\
Indonesia (JCI) & --- & --- & $< 1.96$ & --- & --- & 2012--2025 & No \\
\bottomrule
\end{tabular}
\begin{flushleft}
\small \textit{Notes:} ``c2c'' denotes close-to-close returns (including overnight gaps); ``o2o'' denotes open-to-open returns (implementable portion only). The c2c Sharpe ratios are inflated because the overnight opening gap---which mechanically reflects the previous U.S. session---is captured in c2c returns but cannot be traded upon, as it is realized before the signal can be acted on. All $t$-statistics in the ``$t$-stat (o2o)'' column use Newey-West HAC standard errors with automatic bandwidth selection. The \citet{harvey2016} threshold is $t > 3.0$. Panel B combinations were originally computed using c2c returns; o2o-based combination results would be substantially lower. Transaction costs of 0.3\% per switch are applied throughout.
\end{flushleft}
\end{table}

The distinction between c2c and o2o returns is critical, and its interpretation differs depending on the research question. Using close-to-close returns, the Taiwan 10-day SPY momentum strategy generates a net Sharpe ratio of 1.473 ($t = 3.76$), which would exceed the \citet{harvey2016} threshold. When measured using open-to-open returns---which exclude the overnight gap that a trader cannot capture---the Sharpe ratio drops to 0.87 ($t = 2.22$), below the Harvey threshold. Japan shows a similar pattern (c2c Sharpe 1.306 vs.\ o2o 0.78). Approximately 78\% of the c2c alpha is absorbed by the opening gap.

The c2c results confirm a strong and persistent information transmission channel from U.S. to Asian markets ($r = 0.376$), with six Asia-Pacific markets exceeding the Harvey threshold on a c2c basis (Hong Kong $t = 4.12$, Australia $t = 4.04$, Singapore $t = 4.03$, Korea $t = 3.83$, Taiwan $t = 3.76$, Japan $t = 3.69$). The failure of all U.S.-listed ETF controls confirms that the information is priced with a lag in local markets. Rather than interpreting the c2c--o2o gap as a failure, we view it as evidence of a highly efficient price discovery mechanism: the Asian opening auction---a call auction that aggregates pre-market orders accumulated overnight---incorporates approximately 78\% of overnight U.S. information within the first transaction of the trading day. This finding places the Taiwan opening auction among the most efficient cross-market information aggregation mechanisms documented in the microstructure literature, and has implications for market design in emerging economies where call auction efficiency is an active policy concern.

% K515/K516: Overnight gap decomposition and SPY-conditioned gap analysis
\subsection{Overnight Gap as the Primary Channel}
\label{sec:overnight_gap}

A granular decomposition of 0050.TW returns reveals that the overnight opening gap accounts for 87\% of total close-to-close returns over the 2012--2025 sample, confirming that the bulk of daily price movement is determined before the continuous trading session begins. When conditioned on the previous day's SPY return sign, the gap is strongly asymmetric: following SPY up-days, the average 0050.TW opening gap is $+10.73$ basis points per day ($t = 6.845$), whereas following SPY down-days, the gap averages $-8.91$ bp ($t = -5.23$). The intraday (open-to-close) component shows no significant SPY-conditional pattern ($t = 0.82$), confirming that cross-market information is incorporated almost entirely through the opening auction rather than during continuous trading.

This gap-dominance finding has direct implications for strategy implementability. An ETF-based strategy that attempts to trade on the SPY signal faces a fatal timing problem: by the time the Taiwan market opens and the trader can place an order, the opening gap has already absorbed the signal, imposing an estimated cost of 38.5 bp per round-trip (the average absolute gap). However, a futures-based implementation using TAIEX futures (TX), which trade in a pre-market session beginning at 8:45 (15 minutes before the equity market opens), faces a substantially smaller gap cost of approximately 5 bp, suggesting a potentially viable Sharpe ratio of 0.93 for the futures channel. We leave the full analysis of futures-based implementation---including margin requirements, roll costs, and intraday liquidity---to future work.

\subsection{Source of Information Transmission: Time-Zone Gap, Not Factor Exposure}

A critical question is whether the cross-market predictability reflects genuine information transmission or merely common factor exposure (e.g., global equity beta). We address this with three tests:

\textbf{EWT control.} The iShares MSCI Taiwan ETF (EWT) trades on the New York Stock Exchange during U.S. hours. If the alpha were driven by common factor exposure, the same momentum signal applied to EWT should produce comparable results. It does not: the EWT control fails to achieve statistical significance ($t < 1.96$). This confirms that the alpha originates from the time-zone information gap---the local market's delayed incorporation of U.S. information---rather than from mechanical factor exposure.

\textbf{Europe-to-U.S. test.} If time-zone arbitrage is a universal phenomenon, European close-to-U.S.-open signals should also be profitable. We test whether lagged European returns predict next-day U.S. returns and find a negative correlation ($r = -0.07$), the opposite of the predicted sign. This null result is consistent with our mechanism: the U.S. and European markets have substantial trading-hour overlap (approximately 2.5 hours), which largely eliminates the information gap. Time-zone arbitrage requires \textit{zero} trading overlap.

\textbf{India and Indonesia controls.} India (NSE) and Indonesia (JCI) have partial overlap with U.S. futures trading or limited after-hours information flow. Both fail to exceed the $t > 3.0$ threshold, further supporting the zero-overlap requirement.

\subsection{Robustness}

We subject the time-zone momentum result to five robustness tests. We note that these tests were originally conducted using c2c returns; the statistical significance of the \textit{information transmission channel} is robust, but implementable (o2o) returns are substantially lower as discussed above.

\textbf{Multi-period out-of-sample.} We divide the sample into five non-overlapping sub-periods (approximately three years each). The c2c strategy produces positive Sharpe ratios in all five periods. The combined c2c $t$-statistic across periods is 4.47--4.60, depending on the aggregation method.

\textbf{Bootstrap confidence intervals.} Using 10,000 bootstrap replications (block bootstrap with block length of 20 days to preserve serial dependence), the 95\% confidence interval for the c2c Sharpe ratio is $[0.65, \, 2.24]$, excluding zero.

\textbf{Bonferroni correction.} With six markets tested, the Bonferroni-adjusted significance threshold is $t > 3.3$ (assuming $\alpha = 0.05$ and six independent tests). On a c2c basis, Taiwan ($t = 3.76$), Japan ($t = 3.34$ for the 10-day lookback; $t = 3.69$ for the 5-day lookback), Hong Kong, Australia, Singapore, and South Korea all pass this adjusted threshold. On an o2o basis, none of the individual markets pass, underscoring the gap between information content and implementable alpha.

\textbf{Transaction cost sensitivity.} At the observed 29 switches per year and a conservative 0.3\% per switch, annual transaction costs are approximately 1.7\%. For the c2c strategy the breakeven cost is over 3\% per switch; for the o2o strategy the margin is substantially narrower.

\textbf{Lookback sensitivity.} The 5-day lookback produces a Taiwan c2c Sharpe of 1.141 (NW-adjusted $t = 3.25$); the 10-day lookback produces 1.473 (NW-adjusted $t = 3.76$). Both exceed the \citet{harvey2016} threshold on a c2c basis. The 20-day lookback deteriorates, suggesting that the optimal information horizon is 1--2 trading weeks.

\subsection{Combination Strategies}

The low correlation between the Taiwan and Japan time-zone signals ($\rho = 0.484$) and between the Taiwan time-zone signal and U.S. VT returns ($\rho = 0.14$) suggests potential diversification benefits from combining strategies.

\textbf{TW + JP 50/50.} An equal-weighted combination of Taiwan and Japan time-zone momentum produces a c2c Sharpe ratio of 1.81 ($t = 4.86$), with a maximum drawdown of $-11.9\%$.

\textbf{Global composite (US VT + TW TZ).} Combining a U.S. 12/VIX VT strategy with the Taiwan time-zone strategy produces a c2c Sharpe ratio of 1.61 ($t = 4.31$) and a maximum drawdown of $-8.4\%$.

\textbf{Important caveat.} These combination Sharpe ratios are computed using c2c returns and therefore include the opening-gap component documented above. Since the TZ component's o2o Sharpe is approximately 0.87 (versus c2c 1.47), the implementable combination Sharpe ratios would be substantially lower than reported. The combination results are best interpreted as measuring the \textit{information-theoretic} diversification potential between the VIX-driven variance management channel and the cross-market information transmission channel. The efficient incorporation of overnight U.S. information by the opening auction means that much of the cross-market signal value is already priced in by the time a trader can act.


% ─── Section 6: Macroeconomic Indicators ─────────────────────────

\section{Macroeconomic Indicators for Taiwan Volatility}
\label{sec:macro}

While our three core contributions focus on leverage amplification, VT strategies, and time-zone arbitrage, we briefly examine whether Taiwan-specific macroeconomic indicators provide incremental volatility forecasting power beyond VIX. An important question for practitioners is whether Taiwan-specific macroeconomic variables can enhance volatility forecasts or VT strategies beyond what VIX alone provides. We conduct a comprehensive sweep of 27 Taiwanese macroeconomic indicators within a GARCH-MIDAS framework, testing each variable's incremental contribution to 0050.TW volatility forecasting.

\subsection{Import Growth as a Volatility Predictor}

Among the 27 indicators tested, only one achieves out-of-sample significance: year-over-year import growth (partial $r = 0.214$, OOS improvement $+5.6\%$ over base GJR-GARCH, Diebold-Mariano $p = 0.043$). Taiwan is a highly trade-dependent economy, with imports closely tracking global semiconductor demand and industrial production. When import growth exceeds 20\%---indicating overheated external demand---adding a ``macro guard'' that reduces VT exposure improves the Sharpe ratio by 0.056 but worsens maximum drawdown by 3.2 percentage points, an unfavorable tradeoff.

\subsection{Null Results}

The following indicators fail to improve out-of-sample volatility forecasts or VT performance after controlling for VIX:

\begin{itemize}
    \item \textbf{Business cycle indicator levels}: The National Development Council's composite business cycle score (``traffic light'' system) has no predictive power for next-month realized volatility ($t = -0.53$, $p = 0.60$), consistent with its well-documented lagging nature. Industrial production and M2 money supply growth likewise show insignificant incremental explanatory power (partial $r < 0.05$). However, the \textit{directional change} in these indicators tells a different story; see Section~\ref{sec:bci_mom} below.
    \item \textbf{Margin trading}: Net margin balance and short interest do not predict future volatility or returns. Margin trading in Taiwan appears to be a coincident or lagging indicator rather than a leading indicator.
    \item \textbf{Foreign institutional flows}: Net foreign buying/selling of Taiwanese equities is a lagging indicator, reacting to price movements rather than preceding them.
    \item \textbf{Put/call ratio}: The Taiwan options market put/call ratio shows a contemporaneous correlation of $r = 0.41$ with same-day volatility, but the lagged version drops to zero, confirming that it contains no predictive information.
    \item \textbf{TSMC valuation}: As the dominant constituent of 0050.TW, TSMC's price-earnings ratio might be expected to predict volatility. It does not: the relationship is driven by the AI-related structural break in TSMC's earnings, which invalidated historical PE-volatility patterns.
\end{itemize}

These null results are consistent with our broader finding that VIX serves as a ``sufficient statistic'' for volatility-related information in markets with strong U.S. linkages. For U.S. equities, we have documented that VIX absorbs the predictive content of more than ten alternative indicators (credit spreads, yield curves, VVIX, MOVE, CNN Fear \& Greed Index, AAII sentiment, VIX term structure). For Taiwan, VIX is not perfectly sufficient---import growth and business cycle momentum provide incremental value---but the gap is small enough that the practical benefit of monitoring additional indicators is modest for most investors.

\subsection{Business Cycle Indicator Momentum}
\label{sec:bci_mom}

Beyond trade data, we examine Taiwan's composite business cycle indicators published by the National Development Council. As noted above, the business cycle score itself has no predictive power for next-month realized volatility ($t = -0.53$, $p = 0.60$), consistent with its well-documented lagging nature. However, the month-over-month change in the leading indicator composite achieves the strongest predictive signal among all Taiwan-specific variables ($t = 3.74$, $p < 0.001$, $R^2 = 7.1\%$), surpassing even trade data. A coincident indicator momentum strategy---reducing equity exposure when the coincident composite declines for three consecutive months---yields Sharpe 0.732 (OOS 2018--2024: 1.260), substantially exceeding the $8.63/\text{VIX}$ baseline of 0.334 over the same 2018--2024 sub-period (compared to 0.690 over the full 2016--2026 sample; see Table~\ref{tab:sharpe_reconciliation}).\footnote{The difference between sub-period and full-sample Sharpe ratios for the $8.63/\text{VIX}$ strategy reflects the strategy's underperformance during the prolonged low-volatility bull market of 2020--2021, when the VIX-scaled weight held substantially less than full exposure.} The key insight is that the business cycle indicator's value lies in its \textit{directional change}, not its level.

This finding has broader methodological implications: macroeconomic indicators that appear uninformative in level form may contain substantial predictive content in their first differences. The leading indicator momentum captures regime shifts in the Taiwanese economy---transitions between expansion and contraction---that are orthogonal to the high-frequency volatility information embedded in VIX.

% K512: Dividend ex-date volatility effect
\subsection{Ex-Dividend Volatility and VT Implications}
\label{sec:exdiv}

Taiwan's concentrated ETF dividend calendar---0050.TW and 0056.TW typically go ex-dividend within the same two-week window each July---creates a predictable volatility spike. We examine realized volatility in the five trading days following ex-dividend dates and find an increase of $+32\%$ for 0050.TW and $+69\%$ for the high-dividend 0056.TW relative to the 20-day pre-event baseline. The larger effect for 0056.TW is consistent with its higher dividend yield ($\sim$5--7\% versus $\sim$3\% for 0050.TW) and its retail-investor-heavy shareholder base, which generates more post-ex-date selling pressure.

Despite this pronounced volatility effect, the historical fill rate---the proportion of ex-dividend gaps that are fully recovered within 60 trading days---is 79\% for 0050.TW and 90\% for 0056.TW over the 2012--2025 sample. This high fill rate implies that the ex-dividend volatility spike is transient and mean-reverting, driven by mechanical price adjustment and short-term tax-motivated trading rather than fundamental information.

For VT practitioners, a practical question is whether the ex-dividend volatility spike requires special treatment (e.g., temporarily reducing the target volatility or applying a dividend-adjustment filter). We find that it does not: because VT strategies mechanically reduce exposure when realized volatility rises, the EWMA and GARCH models automatically scale down the position during the post-ex-dividend period without any manual intervention. The ex-dividend effect is therefore self-correcting within the VT framework, and no dividend-specific adjustment is necessary.


% ─── Section 7: VaR and Risk Management ─────────────────────────

\section{VaR and Risk Management}
\label{sec:var}

\subsection{VaR Methodology}

We compute one-day Value-at-Risk at the 1\% level using GJR-GARCH conditional volatility combined with the Student-$t$ distribution (degrees of freedom $= 5$):
\begin{equation}
\text{VaR}_{1\%,t} = \hat{\sigma}_t \cdot q_{0.01}^{t(5)}
\label{eq:var}
\end{equation}
where $q_{0.01}^{t(5)} = -3.365$ is the 1\% quantile of the standardized Student-$t$ distribution with 5 degrees of freedom.

We also evaluate the skewed Student-$t$ distribution, which simultaneously estimates degrees of freedom ($\eta$) and skewness ($\lambda$) via maximum likelihood, adapting to each asset's specific tail behavior. For 0050.TW, the estimated parameters are $\eta = 5.2$ and $\lambda = -0.05$ (near-symmetric with moderate fat tails).

\subsection{Backtest Results}

Over the out-of-sample period 2020--2026 (1,501 trading days), the GJR-GARCH + Student-$t$(5) VaR produces 8 violations (0.5\%), within the Basel III Green Zone threshold of 1.0\% (noting that the traffic-light framework is defined for 250-day windows; we apply the 1\% threshold as a benchmark across our full sample). The Kupiec likelihood ratio test does not reject correct unconditional coverage ($p = 0.12$). For comparison, the Normal distribution produces 30 violations (2.0\%), a borderline result that would trigger the Yellow Zone under Basel III.

The EWMA + Student-$t$(5) combination performs comparably, with 8 violations (0.5\%), demonstrating that the simpler EWMA model is sufficient for regulatory VaR compliance in the Taiwanese market. This is a practical finding: Taiwan investors can implement a fully compliant VaR framework using only EWMA (no GARCH optimization required) plus the Student-$t$ critical value.

\subsection{Skewed Student-$t$ for Taiwan}

The skewed Student-$t$ distribution passes the Kupiec test for all six assets in our broader cross-asset panel (SPY, QQQ, GLD, TLT, EEM, and 0050.TW)---the only distribution to achieve a perfect 6/6 pass rate. The Cornish-Fisher expansion achieves 5/6 (failing for QQQ due to excessive conservatism), while the fixed Student-$t$(5) achieves 4/6 and the Normal distribution achieves only 1/6. For Taiwan specifically, both skewed Student-$t$ and fixed Student-$t$(5) perform well, reflecting the near-symmetric conditional return distribution of 0050.TW.

\subsection{Cornish-Fisher Divergence}

An important caveat for the Taiwanese market is that the Cornish-Fisher (CF) VaR expansion diverges when using a rolling window of $w = 2000$ for 0050.TW. The issue is that kurtosis estimates become extremely large (up to 545 in some windows) due to occasional extreme returns that are not washed out by the long window. With $w = 504$, CF-VaR performs well ($p = 0.88$), but the shorter window introduces persistence bias in the GARCH estimates. Winsorization of extreme kurtosis values resolves the divergence, but this adds an ad hoc element that undermines the method's theoretical appeal. For practical purposes, we recommend the skewed Student-$t$ approach for Taiwan VaR, which is both theoretically principled and computationally stable.


% ─── Section 8: Discussion ──────────────────────────────────────

\section{Discussion}
\label{sec:discussion}

\subsection{Currency Risk}

An important practical consideration for Taiwanese investors considering U.S.-linked strategies is currency risk. The TWD/USD exchange rate introduces additional volatility that reduces the Sharpe ratio of SPY-denominated strategies by approximately 18\% when converted to TWD terms. However, the USD tends to appreciate during equity market crises (a safe-haven effect), providing a partial natural hedge: during the 2022 drawdown, USD appreciation added approximately 15\% to returns measured in TWD.

The 8.63/VIX strategy on 0050.TW (Sharpe 0.69 in TWD, no currency risk) modestly outperforms the 12/VIX strategy on SPY converted to TWD (Sharpe 0.67 including FX effects). Given the additional complexity and cost of currency hedging (estimated at approximately 2\% per annum for TWD/USD forwards), we do not recommend FX hedging for Taiwanese investors implementing VIX-based strategies on local equities.

\subsection{ETFs Versus Individual Stocks for Momentum Strategies}

The time-zone arbitrage strategy performs more consistently on the 0050.TW ETF than on individual Taiwanese stocks. Of the ten individual stocks tested, four pass the \citet{harvey2016} $t > 3.0$ threshold: Yuanta High Dividend ETF (0056.TW, $t = 3.44$), Hon Hai (2317.TW, $t = 3.40$), 0050.TW ($t = 3.25$), and Cathay Financial (2882.TW, $t = 3.06$). However, individual stock results exhibit higher variance and are more susceptible to idiosyncratic events (e.g., Chunghwa Telecom shows a \textit{negative} response to U.S. momentum, $t = -1.17$, reflecting its status as a defensive, dividend-oriented stock that benefits from risk-off flows).

ETFs offer a more stable substrate for the momentum strategy because: (1) diversification smooths idiosyncratic noise, (2) the ETF's larger float reduces market impact costs, and (3) the ETF mechanically reflects the broad market's response to U.S. information, which is the channel the strategy exploits. We therefore recommend 0050.TW (or 0056.TW for income-oriented investors) as the preferred implementation vehicle.

\subsection{Practical Implications}

Our results suggest several practical implications for individual investors in Taiwan, ordered by implementation complexity:

\textbf{VIX Step Rule.} The simplest approach requires no computation: if VIX $< 15$, allocate 100\% to 0050.TW; if VIX $\in [15, 25]$, allocate 70\%; if VIX $> 25$, allocate 40\%. This captures the bulk of VT's drawdown protection with monthly rebalancing.

\textbf{$8.63/\text{VIX}$ strategy.} Compute $w = 8.63 / \text{VIX}$ at each month-end, allocating $w\%$ to 0050.TW and $(1-w)\%$ to a short-term bond fund. Transaction cost: approximately 0.15\% per rebalancing.

\textbf{Time-zone momentum.} Check whether the trailing 10-day SPY cumulative return is positive. If yes, hold 0050.TW; if no, hold cash. This requires daily monitoring and generates approximately 29 switches per year (transaction cost approximately 1.7\% per year). \textit{Caveat:} The c2c Sharpe of 1.47 overstates implementable returns; the o2o Sharpe is 0.87 (see Section~\ref{sec:tz}). The efficient incorporation of overnight U.S. information by Taiwan's opening auction means that most of the signal value is already priced in before trading begins. This strategy is best viewed as an information-transmission indicator rather than a standalone trading strategy.

\textbf{Combined strategy.} Overlaying VIX-based position sizing with momentum switching reduces drawdown, but the combination Sharpe ratios reported in Section~\ref{sec:tz} are based on c2c returns and therefore overstate implementable performance.

\subsection{Comparison to Prior Literature}

Our findings both confirm and extend the existing VT literature. Consistent with \citet{moreira2017}, VT's primary benefit is drawdown reduction rather than Sharpe ratio improvement; the Sharpe improvement $t$-statistic is only 0.33, far below conventional significance levels, but the MDD improvement is statistically significant (bootstrap $p = 0.0004$). This result holds in both U.S. and Taiwanese markets, suggesting it is a universal property of VT strategies.

Our time-zone information transmission finding relates to the broader literature on cross-market momentum. \citet{rapach2013} document that U.S. stock returns predict returns in non-U.S. markets, attributing this to the U.S. market's role as a global information leader. Our contribution is twofold: first, we show that this predictability is statistically significant for \textit{locally traded} Asian equities (but not for U.S.-listed ADRs/ETFs of the same underlying markets), confirming the zero trading-hour overlap condition; second, we establish that the Asian opening auction is a highly efficient mechanism for incorporating overnight U.S. information---approximately 78\% of the predictable return is absorbed by the opening gap, before any trader can act on the signal. This result reframes the standard cross-market predictability narrative: rather than an exploitable anomaly, the c2c--o2o gap provides a quantitative measure of price discovery speed across time zones, demonstrating that the call auction mechanism efficiently aggregates pre-market information. The failure of the European-to-U.S. channel, where trading hours overlap, provides a natural placebo test, and the failure of Taiwan's own U.S.-listed ETF (EWT) confirms that the lag arises from the time-zone gap rather than from factor exposure.

The diversification amplification of the leverage effect extends the work of \citet{ang2002} on asymmetric correlations. While Ang and Chen document that correlations increase during downturns, we show that this phenomenon mechanically amplifies the GJR-GARCH $\gamma$ parameter at the index level relative to the constituent level, and that the amplification ratio varies across markets (4.6$\times$ for the broad TAIEX, 1.45$\times$ for the investable 0050.TW, versus 2.8$\times$ for the U.S.). The gap between TAIEX and 0050.TW ratios highlights the importance of distinguishing between broad market indices and investable vehicles when quantifying the leverage effect.

\subsection{Combining Macro Signals with VIX-Based Strategies}

A natural question is whether Taiwan-specific macro signals can enhance the VIX-based framework. We find that combining VIX scaling with the leading indicator direction---using $k = 10$ (aggressive) when the leading indicator is rising and $k = 6$ (conservative) when declining---significantly improves the pure $8.63/\text{VIX}$ strategy over the full 2012--2026 sample (Diebold-Mariano $p = 0.0005$, Sharpe from 0.999 to 1.151; see Table~\ref{tab:sharpe_reconciliation}). The intuition is straightforward: when the leading indicator signals economic expansion, the VT strategy can afford to take more risk; when contraction looms, it should be more defensive. This conditional approach effectively introduces regime awareness into an otherwise purely volatility-driven framework.

However, this monthly-frequency macro signal cannot improve the daily-frequency SPY momentum strategy (DM $p = 0.82$), confirming that effective signal combination requires matching timescales. The SPY momentum signal already captures high-frequency regime information through daily U.S. market movements, leaving no room for a slower-moving macro indicator to add value. This timescale-matching constraint has practical implications: investors using the time-zone arbitrage strategy should not attempt to overlay macro filters, while investors using the monthly VIX-rebalancing strategy can benefit meaningfully from incorporating leading indicator direction.


\subsection{TSMC Concentration Robustness}
\label{sec:tsmc}

A natural concern is whether our VT results for 0050.TW are driven entirely by its dominant constituent, TSMC, which accounts for approximately 50\% of the index weight and has seen its rolling beta to 0050.TW rise from 0.38 in 2009 to 0.72 by 2026. We address this concern with a decomposition analysis.

We apply VT separately to TSMC (2330.TW) and to a synthetic ex-TSMC 0050.TW portfolio (constructed by removing TSMC's estimated contribution from daily 0050.TW returns). TSMC VT achieves a Sharpe ratio of 1.121, exceeding the 0050.TW VT Sharpe of 0.936. Importantly, the ex-TSMC VT portfolio also produces positive risk-adjusted returns (Sharpe ranging from 0.193 to 0.637 depending on assumptions about TSMC's time-varying weight), confirming that VT benefits are not solely attributable to TSMC.

An unexpected finding from this decomposition concerns the leverage effect. The 0050.TW index exhibits a GJR-GARCH leverage parameter of $\gamma = 0.124$ ($t = 2.46$), whereas TSMC alone shows $\gamma = 0.054$ ($t = 1.07$, insignificant). This result---stronger leverage in the diversified index than in its dominant constituent---reinforces the diversification amplification mechanism documented in Section~\ref{sec:leverage}: aggregation across stocks ``cleans'' the systematic leverage signal by averaging out firm-specific idiosyncratic noise, leaving a purer asymmetric volatility response. TSMC's relatively weak individual leverage effect is consistent with its status as a growth-oriented technology stock whose volatility dynamics are driven more by earnings and semiconductor cycle surprises than by the classical leverage channel of \citet{black1976}.

We acknowledge the growing concentration risk: TSMC explains 52.5\% of 0050.TW return variance over the full sample, and its rolling beta has approximately doubled over the sample period. As TSMC's weight continues to increase, the distinction between 0050.TW VT and single-stock TSMC VT will narrow. Investors concerned about concentration may consider the 0056.TW ETF (which excludes TSMC) as an alternative VT substrate, though at the cost of weaker VIX-based signal quality due to the lower U.S. technology exposure.

\subsection{Sharpe Ratio Reconciliation}

Table~\ref{tab:sharpe_reconciliation} reconciles the Sharpe ratios reported throughout the paper for the $8.63/\text{VIX}$ strategy and its variants. The apparent discrepancies arise from differences in sample period, rebalancing frequency, and strategy specification (pure VIX scaling versus the VIX + leading indicator combination).

\begin{table}[H]
\centering
\caption{Sharpe Ratio Reconciliation for VIX-Based Strategies on 0050.TW}
\label{tab:sharpe_reconciliation}
\begin{tabular}{llcc}
\toprule
Strategy & Period & Sharpe & Table/Section \\
\midrule
$8.63/\text{VIX}$ (monthly, lagged) & 2016--2026 (full) & 0.690 & Table~\ref{tab:vt_results} \\
$8.63/\text{VIX}$ (monthly, lagged) & 2018--2024 (sub-period) & 0.334 & Section~\ref{sec:bci_mom} \\
$8.63/\text{VIX}$ + Leading indicator & 2012--2026 (full) & 0.999 & Section~\ref{sec:discussion} \\
VIX + Leading indicator combo & 2012--2026 (full) & 1.151 & Section~\ref{sec:discussion} \\
\bottomrule
\end{tabular}
\begin{flushleft}
\small \textit{Notes:} All strategies use monthly rebalancing with lagged weights (VIX$_{t-1}$ determines the weight for period $t$) and include transaction costs of 0.585\% round-trip. The 2018--2024 sub-period Sharpe is lower because this window captures the prolonged low-volatility bull market of 2020--2021, during which the VIX-scaled strategy was underinvested relative to buy-and-hold. The VIX + Leading indicator combination uses $k = 10$ when the leading indicator rises and $k = 6$ when it declines; its longer sample (2012--2026) reflects the availability of leading indicator data from 2012.
\end{flushleft}
\end{table}

% ─── Section 9: Conclusion ──────────────────────────────────────

\section{Conclusion}
\label{sec:conclusion}

This paper examines volatility targeting strategies in the Taiwan stock market, contributing three sets of findings.

First, we document that the TAIEX exhibits a strong leverage effect ($\gamma = 0.272$, $t = 3.18$) that is amplified 4.6$\times$ relative to individual stock asymmetries for the broad index ($\sim$900 stocks), though the investable 0050.TW (50 stocks, $\gamma = 0.087$) shows a more modest 1.45$\times$ amplification. Both exceed the U.S. amplification ratio of 2.8$\times$ at the broad-index level. This amplification arises from higher correlation asymmetry among Taiwanese equities during market declines and has direct implications for model selection: index-level strategies should account for asymmetric volatility even when individual constituents show weak asymmetry.

Second, we evaluate EWMA, GARCH, GJR-GARCH, and VIX-proxy VT strategies adapted for the Taiwan market. EWMA VT improves the Sharpe ratio from 0.729 to 0.796 while halving maximum drawdown from $-41.3\%$ to $-18.4\%$. The VIX-proxy strategy ($8.63/\text{VIX}$), calibrated using the VIXTWN-to-VIX ratio of 1.39, achieves comparable drawdown protection with only monthly rebalancing. A key finding is the disconnect between forecast accuracy and portfolio performance: GJR-GARCH does not significantly outperform GARCH in QLIKE (DM $p = 0.86$), yet GJR-based VT produces a meaningfully higher Sharpe ratio (+0.114), suggesting that tail-specific accuracy matters more than average accuracy for VT applications.

Third, we document a persistent information transmission channel from U.S. to Asian equity markets driven by non-overlapping trading hours, and establish that the Asian opening auction is a highly efficient mechanism for cross-market price discovery. A 10-day SPY momentum signal predicts next-day 0050.TW close-to-close returns with $r = 0.376$, and six Asia-Pacific markets exceed the \citet{harvey2016} threshold on a c2c basis. Approximately 78\% of the c2c alpha is absorbed by the overnight opening gap---before any trader can act---so that open-to-open returns yield a Sharpe of only 0.87 ($t = 2.22$). Rather than a negative result, this c2c--o2o gap demonstrates that the call auction mechanism efficiently aggregates pre-market information, incorporating overnight U.S. signals within the first transaction of the trading day. The failure of U.S.-listed ETF controls and the European-to-U.S. channel confirms that zero trading-hour overlap is a necessary condition for observable information transmission lags.

Several limitations warrant mention. First, the VIXTWN data history is limited to 64 months (November 2020 to present), which restricts our ability to validate the VIX-to-VIXTWN ratio across different market regimes. As VIXTWN data accumulates, the calibration should be periodically updated. Second, while the time-zone analysis reveals that the Asian opening auction efficiently incorporates overnight U.S. information, the c2c momentum Sharpe ratio (1.47) substantially overstates implementable returns; the o2o Sharpe of 0.87 confirms that the strategy is better characterized as evidence of efficient cross-market price discovery rather than a tradeable anomaly. Third, our VaR analysis uses a rolling window of 2,000 days, which may be slow to adapt to structural breaks in the Taiwanese market; shorter windows (504 days) improve VaR compliance but introduce GARCH persistence bias. Fourth, we test only a single emerging Asian market in depth; replication in other markets with similar characteristics (e.g., South Korea, where preliminary results are promising) would strengthen the generalizability of our findings. Fifth, TSMC's growing weight in 0050.TW (rolling beta increasing from 0.38 to 0.72 over our sample period) means that our results increasingly reflect the dynamics of a single dominant stock; while the TSMC robustness analysis in Section~\ref{sec:tsmc} confirms that VT benefits extend to the ex-TSMC portfolio, future studies should monitor whether continued concentration erodes the diversification amplification mechanism.

Future research directions include: (1) validating the VIXTWN-based strategy when sufficient data become available (currently 64 months; 252+ months recommended for comprehensive evaluation); (2) extending the time-zone arbitrage analysis to intraday data, which could reveal whether the information gap closes within the first minutes of Taiwan trading; and (3) investigating whether the diversification amplification phenomenon has increased over time as passive ETF investing has grown in Taiwan, potentially strengthening the leverage effect at the index level.
